\subsection{B-空间的结构}

设 B 为一个拓扑空间。

定义 1. --- 拓扑 B-空间（或简称 B-空间）是指一个拓扑空间 X，连同一个从 X 到 B 的连续映射 p。映射 p 称为 B-空间 X 的投影。

设 X、X' 为 B-空间，p、\(p'\) 为它们各自的投影。从 X 到 X' 的 B-态射是指一个从 X 到 X' 的连续映射 f，使得 \(p' \circ f = p\) 。

为方便起见，有时用 \((X,p)\) 表示给拓扑空间 X 装备连续映射 p 之后所得到的 B-空间。

两个 B-态射的复合是一个 B-态射。B-空间的同构也称为 B-同构，它们是作为同胚的那些 B-态射。

这样，如果把集合 X 上的 B-空间结构称为由 X 上的一个拓扑与一个连续映射 \(p: X \rightarrow B\) 所组成的数据，那么就可以取 B-态射作为 B-空间结构的态射 (E, IV, p. 11)。

设 X、\(X'\) 为 B-空间。用 \(\mathcal{C}_{\mathrm{B}}(\mathrm{X};\mathrm{X}')\) 表示从 X 到 \(X'\) 的 B-态射所成的集，用 \(\operatorname{Isom}_{\mathrm{B}}(\mathrm{X};\mathrm{X}')\) 表示从 X 到 \(X'\) 的 B-同构所成的集，用 \(\operatorname{Aut}_{\mathrm{B}}(\mathrm{X})\) 表示 X 的 B-自同构所成的集，也就是从 X 到 X 的 B-同构。

设 X 为一个 B-空间，p 为其投影。若给 B 装备以 \(Id_{B}\) 为投影的 B-空间结构，则集合 \(\mathcal{C}_{\mathrm{B}}(\mathrm{B};\mathrm{X})\) 就是 p 的连续截面 (E, II, p. 18, def. 11) 全体所成的集。

设 X 为一个 B-空间，p 为其投影，b 为 B 的一点。X 的子空间 \(\overline{p}^{1}(b)\) 称为 X 在 b 上的纤维（或 p 在 b 上的纤维），记作 \(X_{b}\) 。为使从 X 到 B-空间 \(X'\) 的连续映射 f 是一个 B-态射，必须且只需 \(f(\mathrm{X}_{b})\) 含于 \(X_{b}'\) ，对一切 \(b \in B\) 。

设 X 为一个 B-空间，p 为其投影。设 f 为从拓扑空间 B' 到 B 中的一个连续映射。连续映射 \(g: B' \to X\) 若满足 \(p \circ g = f\) ，则称为 f 到 X 中的连续提升。换言之，若给 B′ 装备以 f 为投影的 B-空间结构，则 f 到 X 的连续提升就是从 B′ 到 X 的 B-态射。

